% Generated draft: requires mathematical and source review.
\documentclass[11pt]{article}
\usepackage[margin=25mm]{geometry}
\usepackage{fontspec}
\setmainfont{Latin Modern Roman}
\usepackage{amsmath,amssymb,mathtools,unicode-math}
\setmathfont{Latin Modern Math}
\usepackage{enumitem,needspace,xurl,hyperref}
\setlength{\parindent}{0pt}
\setlength{\parskip}{5pt}
\setlength{\emergencystretch}{3em}
\newcommand{\N}{\mathbb{N}}
\newcommand{\Z}{\mathbb{Z}}
\newcommand{\Q}{\mathbb{Q}}
\newcommand{\R}{\mathbb{R}}
\newcommand{\C}{\mathbb{C}}
\newcommand{\F}{\mathbb{F}}
\newcommand{\A}{\mathbb{A}}
\newcommand{\PP}{\mathbb{P}}
\newcommand{\E}{\mathbb{E}}
\newcommand{\eps}{\varepsilon}
\newcommand{\dd}{\,\mathrm{d}}
\begin{document}
\section*{UnsolvedMath ID: 11000112}
\textbf{Source problem code:} catalogue entry
\subsection{Definitions and mathematical statement}
Let \(S\) be a closed connected oriented genus-\(g\) surface, and let \(t_a\) denote the positive Dehn twist about a simple closed curve \(a\). Suppose
\[
 t_{a_1}t_{a_2}\cdots t_{a_n}=1\quad\text{in }\operatorname{Mod}_g,
 \qquad n\geq1.
\]
Set
\[
 G=H_1(S;\Z)/\langle[a_1],\ldots,[a_n]\rangle,
 \qquad b_1(G)=\dim_\Q(G\otimes\Q).
\]
Orientations chosen for the curves do not affect the generated subgroup. The question is whether \(b_1(G)\leq g\), equivalently whether their classes span at least a \(g\)-dimensional subspace of \(H_1(S;\Q)\).

The essential-curve convention is material: if null-homotopic curves are allowed, a single trivial twist satisfies the printed relation but leaves \(b_1(G)=2g\). Thus the usual essential-curve hypothesis, or the exact convention that supplies it, must be recovered. Positivity is also part of the intended factorization, not arbitrary signed twists. No assumption that the curves are distinct or nonseparating is added.

The source says a positive answer would bound the first Betti number of the total space of every genus-\(g\) Lefschetz fibration over \(S^2\) with at least one singular fiber by \(g\), and reports that all examples then known satisfied it. This implication requires the intended essential-vanishing-cycle/relative-minimality convention; the excerpt does not state it.

The continuation separately assumes \(g\geq2\), marks \(P\in S\), defines \(\operatorname{Mod}_g^1\) by one marked puncture, and gives the forgetful exact sequence
\[
 1\longrightarrow\pi_1(S,P)\longrightarrow
 \operatorname{Mod}_g^1\xrightarrow{\phi}\operatorname{Mod}_g\longrightarrow1.
\]
That setup is retained as following background, not used to change the factorization question's unstated genus range.
\subsection{Short English statement}
Must the homology classes in a nonempty positive identity factorization span at least half of surface homology, yielding the stated Lefschetz-fibration Betti bound under the intended essential-cycle conventions?
\subsection{Sources}
\begin{itemize}
\item \emph{Problems on Mapping Class Groups and Related Topics}, edited by Farb, Problem 2.7 and following marked-point setup, supplied PDF page 95. \url{https://www.math.uchicago.edu/~farb/papers/mcgbook.pdf}. Recorded access 2026-07-29; extraction pdf-font-structure. No individual author or separate factorization reference is supplied in the record.
\item Local UnsolvedMath registry, \texttt{data/unsolvedmath/problems.json}, permanent ID \(11000112\), exact source code \texttt{catalogue entry}; supplied OPEN-TRIAGE summary mentions partial positive-factorization results without complete citations.
\end{itemize}
\end{document}
